% Source: https://artofproblemsolving.com/community/c3h584346_2014_mathcounts_state_sprint_18__bank_of_zyzzlvaria
% Stable identity: aops-topic-584346-q18
Think of the digits $1,\dots,5$ as vertices, joined when they differ by at least $2$: $1$ is joined to $3,4,5$; $2$ to $4,5$; $3$ to $1,5$; $4$ to $1,2$; $5$ to $1,2,3$. A code is a path through four different vertices. Count by the omitted digit:
\begin{itemize}
\item omit $3$: on $\{1,2,4,5\}$ the edges are $14,15,24,25$, a $4$-cycle $1\!-\!4\!-\!2\!-\!5\!-\!1$: $4$ starting points $\times$ $2$ directions $=8$ paths;
\item omit $1$: on $\{2,3,4,5\}$ the edges $24,25,35$ form the path $4\!-\!2\!-\!5\!-\!3$: $2$ codes;
\item omit $5$: by symmetry ($d\mapsto6-d$) also $2$;
\item omit $2$: edges $13,14,15,35$ on $\{1,3,4,5\}$: Hamiltonian paths $4\!-\!1\!-\!3\!-\!5$, $4\!-\!1\!-\!5\!-\!3$ and reverses: $4$;
\item omit $4$: by symmetry $4$.
\end{itemize}
The total is $8+2+2+4+4=\boxed{20}$.
